% GSEG DDF — Section 4, concise English working draft.
% Content only. Required packages in DDFmain.tex: longtable, tikz, float.
% TikZ libraries: arrows.meta, positioning, fit.
% Manual section numbers match the existing GSEG TRS/DDF style (secnumdepth=-1).
\section{4. Functional Description}

\subsection{4.1 Functional Architecture}

The Ground Segment (GSEG, CI-02) integrates seven subsystems: Ground Station Network (GSN), Command \& Control (C\&C), Data Processing System (DPS), Early Warning System (EWS), Mission Control Center (MCC), Data Archiving System (DAS), and Ground Infrastructure Subsystem (GIS). Its time-critical event-response chain uses GSN, C\&C, DPS and EWS. MCC supervises operations, owns flight dynamics and mission planning, and supports authorized post-warning actions; DAS retains mission records. GIS provides shared computing, networking, time and site infrastructure. Operator review and archive completion are not prerequisites for the first automatic warning (\texttt{GSEG-REQ-ARC-002}).

GSN is the GSEG communications boundary with the Space Segment (SSEG); C\&C automatically selects two LEO spacecraft using available orbital data and coordinates command execution; DPS owns automated ground event classification, authoritative observation fusion and impact assessment; EWS delivers warnings to authorized recipients. The functional allocation is defined in Section~4.2, while the interaction between these functions is defined in Section~4.3.

\subsection{4.2 Function Tree}

Figure~\ref{fig:gseg_function_tree} shows the first-level decomposition of GSEG functions. Function identifiers (F-xx) are provisional DDF design identifiers, not requirement identifiers.

\begin{figure}[H]
\centering
\begin{tikzpicture}[
  >=Stealth,
  root/.style={draw=HeadingBlue,fill=TableRowLightBlue,rounded corners=2pt,align=center,font=\bfseries\small,text width=7.2cm,minimum height=8mm},
  group/.style={draw=TableBorderBlue,fill=white,rounded corners=2pt,align=center,font=\bfseries\small,text width=5.7cm,minimum height=7mm},
  fn/.style={draw=TableBorderBlue,fill=white,rounded corners=2pt,align=center,font=\small,text width=5.7cm,minimum height=7mm},
  edge/.style={-Stealth,draw=HeadingBlue,thin}
]
\node[root] (root) at (0,0) {F-00 Ground Segment mission support};
\node[group] (critical) at (-3.15,-1.25) {Critical event-response functions};
\node[group] (support) at (3.15,-1.25) {Supervision, retention and assurance};
\node[fn] (f1) at (-3.15,-2.25) {F-01 SSEG data exchange (GSN)};
\node[fn] (f2) at (-3.15,-3.20) {F-02 Observation coordination (C\&C)};
\node[fn] (f3) at (-3.15,-4.15) {F-03 Event assessment (DPS)};
\node[fn] (f4) at (-3.15,-5.10) {F-04 Warning delivery (EWS)};
\node[fn] (f5) at (3.15,-2.25) {F-05 Mission supervision (MCC)};
\node[fn] (f6) at (3.15,-3.20) {F-06 Mission record retention (DAS)};
\node[fn] (f7) at (3.15,-4.15) {F-07 Cross-cutting assurance};
\node[fn] (f8) at (3.15,-5.10) {F-08 Common infrastructure (GIS)};
\draw[edge] (root.south) -- (critical.north);
\draw[edge] (root.south) -- (support.north);
\draw[draw=TableBorderBlue,dashed] (-6.13,-1.76) rectangle (-0.17,-5.57);
\draw[draw=TableBorderBlue,dashed] (0.17,-1.76) rectangle (6.13,-5.57);
\end{tikzpicture}
\caption{GSEG first-level function tree and functional groupings.}
\label{fig:gseg_function_tree}
\end{figure}

The allocation below defines the responsible subsystem at GSEG level; internal decomposition remains subject to the corresponding subsystem DDFs.

{\DDFTableStyle{3}%
\begin{longtable}{|L{0.14}|L{0.72}|L{0.14}|}
\caption{GSEG function allocation}\label{tab:gseg_function_allocation}\\
\hline\rowcolor{TableHeaderBlue}
\thead{ID} & \thead{Function} & \thead{Owner}\\\hline
\endfirsthead
\hline\rowcolor{TableHeaderBlue}\thead{ID} & \thead{Function} & \thead{Owner}\\\hline
\endhead
F-01.1 & Receive and route GEO Event Reports. & GSN\\ \hline
F-01.2 & Transfer LEO Tracking Requests to the SSEG interface. & GSN\\ \hline
F-01.3 & Receive and route LEO Observation Packages and the Onboard Preliminary Trajectory. & GSN\\ \hline
F-01.4 & Exchange spacecraft telemetry and communications status. & GSN\\ \hline
F-02.1 & Evaluate LEO observation opportunities and select two spacecraft. & C\&C\\ \hline
F-02.2 & Generate tracking requests and coordinate their handover. & C\&C\\ \hline
F-02.3 & Track command transmission and execution status. & C\&C\\ \hline
F-03.1 & Associate and validate observation inputs against the event record. & DPS\\ \hline
F-03.2 & Automatically classify valid assessments and determine warning eligibility; distinguish the separate Assessment Unavailable status. & DPS\\ \hline
F-03.3 & Fuse the two selected LEO observation sets and onboard preliminary trajectory. & DPS\\ \hline
F-03.4 & Produce refined trajectory, probabilistic impact area and uncertainty information. & DPS\\ \hline
F-03.5 & Prioritize Baltic-region threats and update assessments as new data arrive. & DPS\\ \hline
F-04.1 & Generate and automatically deliver the first preliminary warning only for warning-eligible outcomes, including the corresponding uncertainty. & EWS\\ \hline
F-04.2 & Track warning delivery status and available acknowledgements. & EWS\\ \hline
F-04.3 & Disseminate authorized post-warning confirmation, revision or cancellation. & EWS\\ \hline
F-05.1 & Present spacecraft and service status, all event-assessment outcomes and warning issue/non-issue decisions and reasons. & MCC\\ \hline
F-05.2 & Support authorized operator actions after first-warning delivery. & MCC\\ \hline
F-05.3 & Coordinate orbital-conjunction information for mission supervision and SSEG. & MCC\\ \hline
F-05.4 & Produce orbit predictions and mission plans for operational use, including current orbital products for C\&C. & MCC\\ \hline
F-06.1 & Record mission products, commands, classification or unavailable-assessment status, warning issue/non-issue decisions and reasons, warning history and audit information. & DAS\\ \hline
F-06.2 & Retain and provide authorized retrieval of mission history. & DAS\\ \hline
F-07.1 & Enforce access control and command authenticity at relevant functions. & Shared\\ \hline
F-07.2 & Protect information and preserve event/product traceability. & Shared\\ \hline
F-07.3 & Detect loss of critical services and support redundant functional paths. & Shared\\ \hline
F-08.1 & Provide shared compute and internal network infrastructure for GSEG services. & GIS\\ \hline
F-08.2 & Provide site power, cooling and common time-reference services. & GIS\\ \hline
F-08.3 & Monitor and maintain availability of common site infrastructure. & GIS\\ \hline
F-08.4 & Maintain shared execution capacity for concurrent critical event workflows. & GIS\\ \hline
\end{longtable}}

MCC owns the operational coordination of orbital-conjunction information (F-05.3). The physical ingress route and forwarding implementation are defined by the relevant interface design. Flight-dynamics products from MCC are made available for C\&C observation planning; their routine preparation and operator review do not gate automatic two-LEO selection after a GEO report. GIS provides enabling infrastructure, while security and mission-product traceability remain shared responsibilities across subsystems. DPS classification does not gate C\&C LEO selection.

\newpage
\subsection{4.3 Functional Chains and Operational States}

\subsubsection{4.3.1 Critical Warning Chain}

The critical chain consists of two GSEG processing intervals separated by SSEG tracking activities (Figure~\ref{fig:gseg_critical_chain}):

\begin{itemize}
\item \textbf{Observation tasking:} GSN receives the GEO report; C\&C selects two LEO spacecraft and prepares tracking requests for GSN handover (2~s maximum; \texttt{GSEG-REQ-PERF-001}).
\item \textbf{Assessment and conditional warning:} after GSN receives both required LEO Observation Packages and the Onboard Preliminary Trajectory, DPS validates the data, classifies the event and determines warning eligibility. For warning-eligible events, EWS automatically delivers the first preliminary warning at the authorized recipient interface (\texttt{GSEG-REQ-PERF-002}; 15~s maximum).
\end{itemize}

\begin{figure}[H]
\centering
\resizebox{0.98\linewidth}{!}{%
\begin{tikzpicture}[
  >=Stealth,
  stage/.style={draw=TableBorderBlue,fill=TableRowLightBlue,rounded corners=3pt,align=center,font=\small,text width=2.25cm,minimum height=1.25cm,inner sep=3pt},
  external/.style={stage,fill=white,text width=2.45cm},
  assess/.style={stage,text width=3.1cm,minimum height=1.50cm},
  outcome/.style={stage,fill=white,text width=3.6cm,minimum height=1.55cm},
  record/.style={stage,text width=3.6cm,minimum height=1.35cm},
  arrow/.style={-{Stealth[length=2mm]},draw=HeadingBlue,line width=0.8pt},
  link/.style={draw=HeadingBlue,line width=0.8pt},
  timing/.style={font=\scriptsize\bfseries,text=HeadingBlue,align=center},
  junction/.style={circle,fill=HeadingBlue,inner sep=1.2pt}
]
\node[external] (geo) at (-6,0) {SSEG GEO\\Event Report};
\node[stage] (gsnup) at (-3,0) {GSN\\report reception};
\node[stage] (cc) at (0,0) {C\&C\\two LEO selected};
\node[stage] (handover) at (3,0) {GSN\\request handover};
\node[external] (sseg) at (6,0) {SSEG LEO\\tracking and return};
\node[stage] (gsndown) at (6,-2.6) {GSN\\LEO receipt /\\product status};
\node[assess] (dps) at (2.3,-2.6) {\textbf{DPS}\\validate both LEO\\classify and assess};
\node[stage] (ews) at (-1.7,-2.6) {\textbf{EWS}\\Automatic\\warning};
\node[external,text width=2.8cm] (recipient) at (-5.5,-2.6) {Authorized\\recipient interface};
\node[outcome] (nonthreat) at (-3.8,-5.4) {\textbf{Credibly non-threatening}\\No external warning};
\node[outcome] (unavailable) at (3.8,-5.4) {\textbf{Assessment Unavailable}\\No external warning};
\node[record] (mcc) at (-3.8,-7.65) {\textbf{MCC}\\Outcome / status\\and reason};
\node[record] (das) at (3.8,-7.65) {\textbf{DAS}\\Decision, reason\\and source-data record};
\draw[arrow] (geo.east) -- (gsnup.west);
\draw[arrow] (gsnup.east) -- (cc.west);
\draw[arrow] (cc.east) -- (handover.west);
\draw[arrow] (handover.east) -- (sseg.west);
\draw[arrow] (sseg.south) -- (gsndown.north);
\draw[arrow] (gsndown.west) -- (dps.east);
\draw[arrow] (dps.west) -- (ews.east);
\node[font=\scriptsize,text=HeadingBlue,fill=white,inner sep=1pt] at (0.09,-2.32) {Eligible};
\draw[arrow] (ews.west) -- (recipient.east);
% DPS supplies one of the two no-warning outcomes; the decision table defines eligibility.
\draw[link] (dps.south) -- (2.3,-4.05);
\draw[link] (-3.8,-4.05) -- (3.8,-4.05);
\node[junction] at (2.3,-4.05) {};
\draw[arrow] (-3.8,-4.05) -- (nonthreat.north);
\draw[arrow] (3.8,-4.05) -- (unavailable.north);
% Both no-warning outcomes are reported to MCC and retained by DAS.
\draw[link] (nonthreat.south) -- (-3.8,-6.65);
\draw[link] (unavailable.south) -- (3.8,-6.65);
\draw[link] (-3.8,-6.65) -- (3.8,-6.65);
\foreach \x in {-3.8,3.8}{\node[junction] at (\x,-6.65) {};}
\draw[arrow] (-3.8,-6.65) -- (mcc.north);
\draw[arrow] (3.8,-6.65) -- (das.north);
\node[timing] at (0,-0.95) {GSN report receipt to request handover: $\leq 2$ s};
\node[timing,text width=6.1cm] at (-3.7,-1.45) {Warning-eligible delivery: $\leq 15$ s\\from complete LEO reception at GSN};
\end{tikzpicture}}
\caption[Conditional critical warning chain and DPS decision branches.]{Conditional critical warning chain. The 2 s interval runs from GSN report reception to tracking-request handover. The 15 s interval runs from complete required two-LEO product reception at GSN to authorized-interface delivery for warning-eligible events; SSEG tracking is outside both GSEG intervals. Both no-warning outcomes are reported to MCC and recorded in DAS. Dots denote connected branches.}
\label{fig:gseg_critical_chain}
\end{figure}

DPS distinguishes three threat classes for valid assessments and a separate data-validity status, \emph{Assessment Unavailable}. Table~\ref{tab:gseg_warning_decision} defines the integrated response; unavailability is not a fourth threat class.

{\DDFTableStyle{2}%
\begin{longtable}{|L{0.47}|L{0.53}|}
\caption{DPS assessment outcomes and conditional GSEG response}\label{tab:gseg_warning_decision}\\\hline
\rowcolor{TableHeaderBlue}\thead{DPS outcome} & \thead{Integrated GSEG response}\\\hline\endfirsthead
\hline\rowcolor{TableHeaderBlue}\thead{DPS outcome} & \thead{Integrated GSEG response}\\\hline\endhead
\texttt{POTENTIAL\_THREAT}\newline Valid two-LEO assessment; refined trajectory and impact-area estimate. & EWS automatically delivers the first preliminary warning to the authorized recipient interface.\\\hline
\texttt{INDETERMINATE\_\allowbreak CREDIBLE\_\allowbreak THREAT\_\allowbreak NOT\_\allowbreak EXCLUDED}\newline Valid assessment; a credible threat cannot be excluded. & EWS automatically delivers the first preliminary warning with explicit uncertainty information.\\\hline
\texttt{CREDIBLY\_\allowbreak NON\_\allowbreak THREATENING}\newline Valid assessment; a credible threat is excluded. & No external warning. MCC receives the result and reason; DAS retains the decision, reason and source data.\\\hline
\texttt{ASSESSMENT\_\allowbreak UNAVAILABLE}\newline Missing, invalid or insufficient required observations or products. & No external warning. Report inability to assess to MCC; retain the status and reason in DAS.\\\hline
\end{longtable}}

Under \texttt{GSEG-REQ-FUN-005}, the first warning requires valid assessment of both selected LEO observation sets, a refined ground trajectory and an impact-area estimate, and a potential threat or an indeterminate assessment for which a credible threat cannot be excluded. DPS determines the classification and warning eligibility; EWS distributes the warning. MCC review and DAS write completion do not gate the first warning. Thresholds distinguishing credible threats and non-threatening events remain TBD for DPS/EWS design and DDJF substantiation; no numerical decision threshold is established here.

\subsubsection{4.3.2 Supporting and Update Chains}

C\&C, GSN, DPS and EWS provide relevant operational and event status to MCC, including all classification and warning-decision outcomes. DAS records applicable mission products, commands, warning issue/non-issue decisions and reasons, warning versions and audit information in parallel with the critical chain. New observations trigger updated DPS assessments; an authorized operator may subsequently confirm, revise or cancel an issued warning through MCC and EWS.

\subsubsection{4.3.3 Operational States and Degraded Operation}

\begin{itemize}
\setlength{\itemsep}{0pt}
\setlength{\parsep}{0pt}
\setlength{\parskip}{0pt}
\item \textbf{Nominal monitoring:} constellation and ground-service supervision, with event-response functions ready.
\item \textbf{Event response:} GEO report reception, LEO selection and tracking-request handover.
\item \textbf{Event assessment:} receipt and validation of LEO products, classification and ground fusion.
\item \textbf{Warning dissemination, for eligible events:} automatic delivery of the first preliminary warning.
\item \textbf{Post-warning operations:} updated assessments, operator actions and warning updates.
\end{itemize}

These are event stages, not mutually exclusive GSEG modes. Assessment Unavailable produces no external warning and status/records for MCC/DAS; service loss is reported to MCC. Ten concurrent events are a proposed sizing case (P), not a mission maximum. MCC prepares orbital products in advance; C\&C selection does not wait for DPS classification, fresh flight-dynamics computation or operator approval. Physical recovery is defined outside Section~4.
